% answers_short.tex —— 2003 年·北京卷·理科 · 简短答案速查
% 由 tools/make_short_answers.py 自 answers.tex 的 \gksolhead/\gkanswer 提取，
% 并与 2003/2003_数学真题_北京理_答案要点.md 逐题比对后生成，勿手改（改源 MD 后重跑）。

\begin{gknotice}
\textbf{使用说明}：本节只给最终答案便于快速核对，完整推导见解析版。答案为编写者独立推导，\textbf{非官方标准答案}。
\end{gknotice}

\gkpart{选择题}
\begin{center}\begin{tabular}{ccccccccccc}
\toprule
题号 & 1 & 2 & 3 & 4 & 5 & 6 & 7 & 8 & 9 & 10 \\
\midrule
答案 & A & D & A & B & D & B & C & B & C & C \\
\bottomrule
\end{tabular}\end{center}

\gkpart{填空题}
\begin{itemize}[leftmargin=2.6em]
  \item[\textbf{11.}] \(f(x)\) 与 \(g(x)\)
  \item[\textbf{12.}] \(y^2=-36(x-4)\)
  \item[\textbf{13.}] \(V=\dfrac{\pi r^2(a+b)}2\)
  \item[\textbf{14.}] \(\dfrac4{4+\pi}\)
\end{itemize}

\gkpart{解答题}
\begin{itemize}[leftmargin=2.6em]
  \item[\textbf{15.}] (1) 最小正周期为 \(\pi\)；(2) 最大值为 1，最小值为 \(-\sqrt2\)
  \item[\textbf{16.}] (1) \(a_n=2n\)；(2) \(x=1\) 时 \(S_n=n(n+1)\)；\(x\ne1\) 时 \(S_n=\dfrac{2x(1-(n+1)x^n+nx^{n+1})}{(1-x)^2}\)
  \item[\textbf{17.}] (1) \(BC_1\parallel\) 平面 \(AB_1D\)；(2) U/PENDING；若另有证据确认底面边长 \(s=3\)，条件性结果为 \(60^\circ\)；(3) U/PENDING；若另有证据确认底面边长 \(s=3\)，条件性结果为 \(V=27/8\)
  \item[\textbf{18.}] (1) \(\dfrac{x^2}{a^2}+\dfrac{(y-r)^2}{b^2}=1\)，焦点为 \((\pm\sqrt{a^2-b^2},r)\)，离心率为 \(\dfrac{\sqrt{a^2-b^2}}a\)；(2) \(\dfrac{kx_1x_2}{x_1+x_2}=\dfrac{r^2-b^2}{2r}\)；(3) \(|OP|=|OQ|\)
  \item[\textbf{19.}] (1) \(P=(0,h/3)\)，其中 \(h=\sqrt{a^2-b^2}\)；(2) \(h\le b\) 时 \(P=O\)；\(h>b\) 时 \(P=(0,\dfrac{h^2-b^2}{2h})\)
  \item[\textbf{20.}] (1) \(x-1\le f(x)\le1-x\)；(2) \(|f(u)-f(v)|\le1\)；(3) 不存在满足条件的奇函数
\end{itemize}
